\chapter{语言模型}
\section{语言模型简介}
回忆一下上文讲的 CBOW 模型，假设单词序列为$w_1, w_2, ··· , w_T$且窗口大小为1，
它所做的事情就是从上下文$w_{t−1}$ 和 $w_{t+1}$预测目标词$w_t$。
当给定上下文时，目标词是 $w_t$ 的概率为$P(w_t \mid w_{t-1}, w_{t+1})$，CBOW 模型就是对这一后验概率进行建模，
如果使用交叉熵误差，其损失函数就是$L = -\log P(w_t \mid w_{t-2}, w_{t-1})$。在上一章我们使用 CBOW 模型的权重获得单词的表示，
那么它本来的目标“从上下文预测目标词”是否能用来做点事情呢？这就得提到语言模型了。

语言模型（language model）是一种对句子”打分“的模型——它判断一个句子是否自然、通顺。具体来说，语言模型输出概率来衡量一个单词序列的合理性。
比如，对于“you say goodbye”这一单词序列，语言模型给出高概率；
对于“you say good die”这一单词序列，模型则给出低概率。
语言模型常用于机器翻译和语音识别等任务，比如，语音识别系统根据人的发言生成多个句子作为候选，
此时，使用语言模型，可以按照单词序列发生的可能性对候选句子进行排序。

困惑度（perplexity）常被用作评价语言模型的预测性能的指标，困惑度表示“概率的倒数”。
假设“you say goodbye and i say hello.”是语料库中一个句子，而且在向语言模型传入单词 you 时会输出的概率分布中，
单词是 say 的概率为 0.2，那么困惑度就是 1/0.2=5。困惑度可以解释为语言模型做下一个选择的选项的数量，
在刚才的例子中，模型的困惑度是 5，这意味着下一个要出现的单词的候选个数为 5 个。
在输出数据为多个的情况下，困惑度计算公式为：$Perplexity = e^L$，其中，L为交叉熵损失函数的值。

用数学表示语言模型：
\begin{align*}
	P(w_1, \dots, w_m) & = P(w_m \mid w_1, \dots, w_{m-1}) \cdot P(w_{m-1} \mid w_1, \dots, w_{m-2}) \cdots P(w_2 \mid w_1) \cdot P(w_1) \\
	                   & = \prod_{t=1}^{m} P(w_t \mid w_1, \dots, w_{t-1})
\end{align*}

假设我们将CBOW模型改造成语言模型，那么首先得将原来对称的上下文改成只有目标词左侧的上下文，这样可以近似实现语言模型：
$$
	P(w_1, \dots, w_m) = \prod_{t=1}^{m} P(w_t \mid w_1, \dots, w_{t-1}) \approx \prod_{t=1}^{m} P(w_t \mid w_{t-2}, w_{t-1})
$$


\section{RNN}
语言模型虽然能用 CBOW 来近似实现，但存在两个明显的局限。
第一，它的上下文窗口大小是固定的，只能看到有限个过去的单词，更远的信息会被丢弃。
第二，它对上下文中的单词做平均，因此单词的顺序信息会被忽视，比如，(you, say) 和 (say, you) 会被作为相同的内容进行处理。
循环神经网络（Recurrent Neural Network，RNN）可以很好地解决上述问题。

之前我们看到的神经网络都是前馈型神经网络，网络的传播方向是单向的，没有环路，
虽然前馈网络结构简单、易于理解，但它不能很好地处理时间序列数据。
RNN 与前馈网络的关键区别在于它引入了环路，使得网络能够维持一个内部状态，从而捕捉序列中的时序依赖关系。
\begin{figure}[H]
	\centering
	\includegraphics[width=0.35\textwidth, keepaspectratio]{5-1.pdf}
	\caption{RNN层示例}
\end{figure}

RNN 通过环路，使得数据可以在层内循环。时序数据 $(x_0, x_1, ··· , x_t, ···)$ 被输入到层中。
然后，以与输入对应的形式，输出 $(h_0, h_1, ··· , h_t, ···)$。RNN 有两个权重，
分别是转化输入$x$的权重$W_x$和转化前一个 RNN 层的输出的权重$W_h$。
RNN的计算用数学表示为：$h_t = \tanh(h_{t-1} W_h + x_t W_x + b)$。
RNN 的输出$h_t$称为隐藏状态（hidden state）或隐藏状态向量（hidden state vector）。

由于RNN 层有环路，因此不能直接使用前馈神经网络的反向传播算法，
RNN的反向传播被称为基于时间的反向传播（Back Propagation Through Time，BPTT），
具体做法是展开RNN的循环，把它看作普通的神经网络，然后运用之前的反向传播方法即可。
\begin{figure}[H]
	\centering
	\includegraphics[width=0.9\textwidth, keepaspectratio]{5-2.pdf}
	\caption{RNN的展开}
\end{figure}

在处理长度为 1000 的时序数据时，如果展开 RNN 层，它将成为在水
平方向上排列有 1000 个层的网络，序列太长的话，会出现计算量或者内存使用量方面的问题，
此外，随着层变长，梯度逐渐变小，最终将无法传递下去。
因此，考虑将时间轴方向上过长的网络在合适的位置进行截断，从而创建多个小型网络，
然后对截出来的小型网络各自执行误差反向传播法，这个方法称为 Truncated BPTT。

在RNN的学习中，为了执行 mini-batch 学习，需要进将输入数据进行“偏移”。
对长度为1000的时序数据，如果将批大小设为 2，那么第 1 笔样本数据从头开始按顺序输入，
第 2 笔数据就从第 500 个数据开始按顺序输入。

利用RNN可以构建基于神经网络的语言模型，称为RNN语言模型（RNN Language Model，RNNLM）。
RNNLM的结构很简单：输入层接收单词的向量表示（如词嵌入），经过RNN层处理后得到隐藏状态，
最后通过全连接层和softmax输出下一个单词的概率分布。

\begin{figure}[H]
	\centering
	\includegraphics[width=0.2\textwidth, keepaspectratio]{5-3.pdf}
	\caption{RNNLM}
\end{figure}

\section{LSTM和GRU}
前面介绍的简单的RNN在无法很好地学习到时序数据的长期依赖关系。RNN中使用了函数$\tanh(x)$，
它的导数为$1 - \tanh^2(x)$，这个值小于1，并且随着 x 远离 0，它的值在变小，
这就导致RNN在处理长序列时容易出现梯度消失的问题。另外，RNN中使用了矩阵乘积，
有可能导致梯度消失或梯度爆炸，这里不具体分析了。

梯度爆炸的问题比较好解决，可以通过梯度裁剪（gradients clipping）来解决，
$$
	g \leftarrow
	\begin{cases}
		g,                             & \text{if } \|g\|_2 \leq \tau \\
		g \cdot \dfrac{\tau}{\|g\|_2}, & \text{if } \|g\|_2 > \tau
	\end{cases}
$$
这里假设可以将神经网络用到的所有参数的梯度整合成一个向量 $g$，且阈值设置为 $\tau$。

梯度消失的问题则更为棘手，为了解决这个问题，LSTM（Long Short-Term Memory）引入了一种名为“门”的结构，
并且保存专门的权重参数用于控制门的开合程度，这些权重参数也会在学习中被更新。

在 LSTM 中存在记忆单元 c，仅在 LSTM 层内部接收和传递数据，不向其他层输出。
其中$h_t = \tanh(c_t)$。
\begin{figure}[H]
	\centering
	\includegraphics[width=0.4\textwidth, keepaspectratio]{5-4.pdf}
	\caption{RNN与LSTM}
\end{figure}

输出门（output gate）针对 $\tanh(c_t)$ 的元素，调整它们作为下一时刻的隐藏状态的重要程度。
输出门的开合程度根据输入 $x_t$ 和上一个状态 $h_{t−1}$ 求出($\sigma$表示sigmod函数)，
最后将这个 $o_t$ 和 $\tanh(c_t)$ 的对应元素的乘积作为 $h_t$ 输出。：
$$
	o_t = \sigma (x_t W_x^{(o)} + h_{t-1} W_h^{(o)} + b^{(o)})
$$

遗忘门（forget gate）是一个忘记不必要记忆的门，遗忘门的输出 $f_t$ 与记忆单元 $c_{t-1}$ 相乘，生成新的记忆单元 $c_t$：
$$f_t = \sigma (x_t W_x^{(f)} + h_{t-1} W_h^{(f)} + b^{(f)})$$

遗忘门用来遗忘记忆，现在我们还想向记忆单元添加一些应当记住的新信息，为此我们添加新的 $\tanh$ 节点，
通过将这个 $g_t$ 加到上一时刻的 $c_{t−1}$ 上，从而形成新的记忆:
$$g_t = \tanh ( x_t W_x^{(g)} + h_{t-1} W_h^{(g)} + b^{(g)})$$

最后添加一个输入门（input gate），将 $i_t$ 和 $g_t$ 的对应元素的乘积添加到记忆单元中：
$$i_t = \sigma (x_t W_x^{(i)} + h_{t-1} W_h^{(i)} + b^{(i)})$$

在LSTM中，记忆单元的更新仅有+和x运算，因此梯度消失和梯度爆炸问题都得到有效缓解。
\begin{figure}[H]
	\centering
	\includegraphics[width=0.6\textwidth, keepaspectratio]{5-5.png}
	\caption{LSTM}
\end{figure}

GRU（Gated Recurrent Unit）是LSTM的一种简化变体，同样通过门控机制来解决梯度消失问题。
与LSTM不同，GRU没有记忆单元$c$，只有一个隐藏状态$h$在时间方向上传播。
GRU使用两个门（LSTM使用三个门）：reset门$r$和update门$z$。

reset门决定在多大程度上“忽略”过去的隐藏状态：
$$r_t = \sigma(x_t W_x^{(r)} + h_{t-1} W_h^{(r)} + b^{(r)})$$

update门扮演了LSTM中forget门和input门两个角色，控制隐藏状态的更新：
$$z_t = \sigma(x_t W_x^{(z)} + h_{t-1} W_h^{(z)} + b^{(z)})$$

基于reset门，计算候选隐藏状态$\tilde{h}_t$：
$$\tilde{h}_t = \tanh(x_t W_x + (r_t \odot h_{t-1}) W_h + b)$$
如果$r_t$为0，则候选隐藏状态仅取决于当前输入$x_t$，过去的隐藏状态被完全忽略。

最终，GRU通过update门对过去的隐藏状态和候选隐藏状态进行加权，得到新的隐藏状态：
$$h_t = (1 - z_t) \odot h_{t-1} + z_t \odot \tilde{h}_t$$
其中$(1 - z_t) \odot h_{t-1}$部分充当forget门的功能，从过去的隐藏状态中删除应该被遗忘的信息；
$z_t \odot \tilde{h}_t$部分充当input门的功能，对新增的信息进行加权。

\begin{figure}[H]
	\centering
	\includegraphics[width=0.6\textwidth, keepaspectratio]{5-6.png}
	\caption{GRU}
\end{figure}
